\par\noindent\textbf{(1)} 局域回旋频率 \(\omega_c=|q|B/m\)，垂直能量 \(E_\perp=mv_\perp^2/2\)。快变量作用 \(J_\perp=(2\pi)^{-1}\oint p\,dq\) 与 \(E_\perp/\omega_c\) 成正比，故 \(\mu_{\mathrm{ad}}=E_\perp/B=mv_\perp^2/(2B)\) 在慢变磁场中近似守恒。
\par\noindent\textbf{(2)} 无电场时总动能 \(E=mv_\parallel^2/2+\mu_{\mathrm{ad}}B(s)\) 近似守恒，因此 \(v_\parallel^2(s)=2[E-\mu_{\mathrm{ad}}B(s)]/m\)。磁强增大时，垂直能量上升、平行动能下降，是能量在两方向间交换而非磁场作净功。
\par\noindent\textbf{(3)} 反射点满足 \(v_\parallel=0\)，即 \(B_m=E/\mu_{\mathrm{ad}}\)。在最小磁场 \(B_0\) 处定义俯仰角 \(\sin^2\alpha_0=E_{\perp0}/E=\mu_{\mathrm{ad}}B_0/E\)，故 \(B_m=B_0/\sin^2\alpha_0\)。若 \(B_{\max}\) 是镜区最大场，反射需 \(\sin^2\alpha_0\ge B_0/B_{\max}\)。
\par\noindent\textbf{(4)} 一阶磁矩守恒要求 \(\rho_L/L_B\ll1\)、\(\omega_c^{-1}/T_B\ll1\)，其中 \(L_B=|B/\nabla_\parallel B|\)。靠近磁零点 \(B\to0\) 时回旋频率衰减、Larmor半径增大，粒子可能非绝热散射，磁镜公式无法继续简单外推。
